% ======================================================================
\chapter{Assinatura de commits com GPG e SSH}
\label{cap:assinatura}

\section{O problema: o autor de um commit é só um texto}

No capítulo~\ref{cap:git} configuramos o nome e o e-mail com \comando{git config}. O Git aceita \textbf{qualquer} valor ali, sem conferir nada. Veja:

\begin{bashcode}
git -c user.name="Linus Torvalds" -c user.email="torvalds@linux-foundation.org" \
    commit --allow-empty -m "Commit que parece do Linus"
git log --format='%an <%ae>' -1
\end{bashcode}
\begin{saida}
Linus Torvalds <torvalds@linux-foundation.org>
\end{saida}

O commit foi feito por nós, mas o histórico diz que foi o criador do Git. Se esse commit for enviado ao GitHub, ele aparece com o nome e a foto de quem tem aquele e-mail cadastrado.

\begin{conceito}[Assinatura de commit]
Assinar um commit é anexar a ele uma \textbf{assinatura criptográfica} feita com uma chave privada que só você tem. Quem tem a sua chave \textbf{pública} consegue conferir que o commit veio de você e que não foi alterado depois. O GitHub faz essa conferência e mostra um selo ao lado do commit.
\end{conceito}

\subsection*{Os selos}

\begin{center}
\renewcommand{\arraystretch}{1.35}
\begin{tabularx}{\linewidth}{@{}>{\bfseries}l X@{}}
\toprule
\textbf{Selo} & \textbf{Significado} \\
\midrule
\color{verdeok}Verified & o commit está assinado e a assinatura foi conferida com uma chave cadastrada na conta \\
\color{vermelhoerro}Unverified & o commit está assinado, mas a assinatura não pôde ser conferida (chave não cadastrada, e-mail diferente, chave de outra pessoa) \\
\textrm{(sem selo)} & o commit não foi assinado \\
\bottomrule
\end{tabularx}
\end{center}

\begin{aviso}[O e-mail precisa bater]
Para o selo \textbf{Verified}, o e-mail do commit (\comando{git config user.email}) precisa ser um e-mail \textbf{verificado} na sua conta do GitHub. No caso do GPG, ele também precisa estar dentro da chave. Essa é a causa mais comum de \enquote{Unverified}.
\end{aviso}

\section{GPG ou SSH?}

Há dois jeitos de assinar, e os dois são aceitos pelo GitHub:

\begin{center}
\renewcommand{\arraystretch}{1.35}
\small
\begin{tabularx}{\linewidth}{@{}>{\bfseries\color{roxoescuro}}l >{\raggedright\arraybackslash}X >{\raggedright\arraybackslash}X@{}}
\toprule
 & \textbf{GPG} & \textbf{SSH} \\
\midrule
Ferramenta & GnuPG (\comando{gpg}) & OpenSSH (\comando{ssh-keygen}), que já vem com o Git \\
Requisito & GnuPG instalado & Git 2.34 ou mais novo \\
Chave & uma chave própria, que carrega nome e e-mail & a mesma chave SSH do capítulo~\ref{cap:git} pode servir \\
Validade & pode ter data de expiração embutida & sem data de expiração \\
Conferir no terminal & direto, com o chaveiro do GPG & precisa de um arquivo \arq{allowed\_signers} \\
Quando escolher & padrão consolidado, usado há muito tempo & mais simples para quem já usa SSH \\
\bottomrule
\end{tabularx}
\end{center}

\begin{dica}
Escolha \textbf{um} dos dois e siga só a seção correspondente. Se não souber qual, comece pelo SSH: são menos passos.
\end{dica}

% -----------------------------------------------------------------------------
\section{Assinando com GPG}

\subsection{Instalar o GnuPG}

\begin{itemize}
  \item \textbf{Linux:} costuma vir instalado. Confira com \comando{gpg --version}. No Debian e no Ubuntu, se faltar: \comando{sudo apt install gnupg}.
  \item \textbf{Windows:} o Git for Windows já traz um \comando{gpg} dentro do \textbf{Git Bash}. Use o Git Bash para todos os comandos desta seção.
\end{itemize}

\begin{aviso}[Windows: um GPG só]
Se você também instalar o Gpg4win, passa a ter \textbf{dois} GPGs no computador, cada um com o próprio chaveiro. Uma chave criada num não aparece no outro, e o Git reclama de \enquote{No secret key}. Crie a chave e assine sempre pelo mesmo. Para saber qual o Git está usando, rode \comando{which gpg} no Git Bash.
\end{aviso}

\subsection{Passo 1 — Ver se você já tem uma chave}

\begin{bashcode}
gpg --list-secret-keys --keyid-format=long
\end{bashcode}

Se aparecer uma chave com o seu e-mail, pule para o passo 3.

\subsection{Passo 2 — Gerar a chave}

\begin{bashcode}
gpg --full-generate-key
\end{bashcode}

O comando faz algumas perguntas:

\begin{center}
\renewcommand{\arraystretch}{1.35}
\small
\begin{tabularx}{\linewidth}{@{}>{\bfseries}l >{\raggedright\arraybackslash}X@{}}
\toprule
\textbf{Pergunta} & \textbf{O que responder} \\
\midrule
Tipo de chave & aperte \tecla{Enter} para aceitar o padrão. Nas versões novas é ECC (ed25519); nas antigas, \enquote{RSA e RSA}. O GitHub aceita os dois \\
Tamanho & só aparece para RSA: digite \texttt{4096} \\
Validade & \texttt{0} = não expira, que é o padrão recomendado pelo GitHub. Uma data (como \texttt{1y}, um ano) é mais segura, mas exige renovar a chave \\
Nome e e-mail & o seu nome e \textbf{o mesmo e-mail verificado} no GitHub \\
Senha & uma senha forte. Ela protege a chave se alguém copiar o seu computador \\
\bottomrule
\end{tabularx}
\end{center}

\subsection{Passo 3 — Descobrir o ID da chave}

\begin{bashcode}
gpg --list-secret-keys --keyid-format=long
\end{bashcode}
\begin{saida}
sec   rsa4096/30F2B65B9246B6CA 2017-08-18 [SC]
      D5E4F29F3275DC0CDA8FFC8730F2B65B9246B6CA
uid                   [ultimate] Maria Exemplo <maria@exemplo.com>
ssb   rsa4096/B7ABC0813E4028C0 2017-08-18 [E]
\end{saida}

O ID é o que vem \textbf{depois da barra} na linha \arq{sec}: aqui, \arq{30F2B65B9246B6CA}. Numa chave ECC, a linha começa com \arq{ed25519/}. O seu ID será outro.

\subsection{Passo 4 — Exportar a chave pública}

\begin{bashcode}
gpg --armor --export 30F2B65B9246B6CA
\end{bashcode}

Copie \textbf{tudo}, da linha \arq{-----BEGIN PGP PUBLIC KEY BLOCK-----} até a linha \arq{-----END PGP PUBLIC KEY BLOCK-----}, incluindo as duas.

\begin{aviso}[Pública, nunca privada]
O que vai para o site é a chave \textbf{pública} (\comando{--export}). Nunca use \comando{--export-secret-keys} para isso. A chave privada não sai do seu computador.
\end{aviso}

\subsection{Passo 5 — Cadastrar no GitHub}

\begin{enumerate}
  \item Clique na sua foto (canto superior direito) \(\rightarrow\)\ \textbf{Settings}.
  \item Na seção \textbf{Access}, clique em \textbf{SSH and GPG keys}.
  \item Clique em \textbf{New GPG key}.
  \item Em \textbf{Title}, dê um nome (por exemplo, \enquote{Notebook pessoal}). Em \textbf{Key}, cole o bloco copiado.
  \item Clique em \textbf{Add GPG key}. O GitHub pode pedir a sua senha de novo.
\end{enumerate}

\subsection{Passo 6 — Dizer ao Git qual chave usar}

\begin{bashcode}
git config --global user.signingkey 30F2B65B9246B6CA
git config --global commit.gpgsign true     # assina todo commit
git config --global tag.gpgSign true        # assina toda tag
\end{bashcode}

\begin{dica}[Já usou assinatura SSH antes?]
Se algum dia configurou \comando{gpg.format ssh}, desfaça com \comando{git config --global --unset gpg.format}. Sem isso, o Git continua tentando assinar com SSH.
\end{dica}

\subsection{Passo 7 — Ensinar o GPG a pedir a senha no terminal}

No Linux e no Git Bash, o GPG precisa saber em qual terminal pedir a senha da chave. Rode uma vez:

\begin{bashcode}
[ -f ~/.bashrc ] && echo -e '\nexport GPG_TTY=$(tty)' >> ~/.bashrc
source ~/.bashrc
\end{bashcode}

Quem usa \textbf{zsh} coloca a linha \arq{export GPG\_TTY=\$(tty)} no \arq{\textasciitilde/.zshrc}.

\subsection{Passo 8 — Testar}

\begin{bashcode}
git commit --allow-empty -m "Testando a assinatura"
git log --show-signature -1
\end{bashcode}
\begin{saida}
commit e55fbf5be5b86eb89737adea712d8c605b5f9529
gpg: Assinatura feita seg 28 set 2026 14:51:16 -03
gpg:                usando RSA chave 0D9FFCDB3C4EBDC3E7225C8CB4BEA86F1F6C7956
gpg: Assinatura correta de "Maria Exemplo <maria@exemplo.com>" [final]
Author: Maria Exemplo <maria@exemplo.com>
\end{saida}

A linha \enquote{Assinatura correta} (em inglês, \enquote{Good signature}) confirma que funcionou. Se você não ligou o \comando{commit.gpgsign}, assine um commit específico com \comando{git commit -S -m "..."}.

% -----------------------------------------------------------------------------
\section{Assinando com SSH}

\subsection{Requisitos}

\begin{itemize}
  \item \textbf{Git 2.34 ou mais novo} (\comando{git --version});
  \item \textbf{OpenSSH 8.1 ou mais novo} (\comando{ssh -V}). A versão 8.7 tem a assinatura quebrada: nela, atualize para a 8.8 ou mais nova;
\end{itemize}

\subsection{Passo 1 — Ter uma chave SSH}

Se você fez o login por SSH no capítulo~\ref{cap:git}, já tem uma: normalmente \arq{\textasciitilde/.ssh/id\_ed25519.pub}. Se não tiver, crie:

\begin{bashcode}
ssh-keygen -t ed25519 -C "maria@exemplo.com"
\end{bashcode}

Aperte \tecla{Enter} para aceitar o local padrão e escolha uma senha. O arquivo terminado em \arq{.pub} é a chave \textbf{pública}. O outro, sem extensão, é a \textbf{privada}, e nunca sai do seu computador.

\subsection{Passo 2 — Dizer ao Git para assinar com SSH}

\begin{bashcode}
git config --global gpg.format ssh
git config --global user.signingkey ~/.ssh/id_ed25519.pub
git config --global commit.gpgsign true
git config --global tag.gpgSign true
\end{bashcode}

\subsection{Passo 3 — Cadastrar como chave de assinatura no GitHub}

\begin{enumerate}
  \item Foto \(\rightarrow\)\ \textbf{Settings} \(\rightarrow\)\ \textbf{SSH and GPG keys} \(\rightarrow\)\ \textbf{New SSH key}.
  \item Em \textbf{Title}, dê um nome. Em \textbf{Key type}, escolha \textbf{Signing Key}.
  \item Em \textbf{Key}, cole o conteúdo do arquivo \arq{.pub} (\comando{cat \textasciitilde/.ssh/id\_ed25519.pub}) e clique em \textbf{Add SSH key}.
\end{enumerate}

\begin{aviso}[No GitHub, a mesma chave entra duas vezes]
Uma chave cadastrada como \emph{Authentication Key} serve para \comando{git push}, mas \textbf{não} para verificar assinaturas. Se quiser usar a mesma chave para as duas coisas, cadastre-a duas vezes: uma com cada tipo.
\end{aviso}

Pelo terminal, com o GitHub CLI, dá no mesmo. A opção \comando{--type} existe nas versões do \comando{gh} lançadas a partir de 2023. Se a sua recusar a opção, use a página web.

\begin{bashcode}
gh ssh-key add ~/.ssh/id_ed25519.pub --type signing --title "Notebook pessoal"
\end{bashcode}

\subsection{Passo 4 — Permitir a conferência no seu computador}

O GitHub confere a assinatura sozinho. Mas, para o \emph{seu} Git conferir, ele precisa de uma lista de chaves confiáveis. Sem ela, aparece:

\begin{saida}
error: gpg.ssh.allowedSignersFile needs to be configured and exist for ssh signature verification
\end{saida}

Crie a lista uma vez:

\begin{bashcode}
touch ~/.ssh/allowed_signers
git config --global gpg.ssh.allowedSignersFile ~/.ssh/allowed_signers
echo "$(git config --get user.email) namespaces=\"git\" $(cat ~/.ssh/id_ed25519.pub)" >> ~/.ssh/allowed_signers
\end{bashcode}

Cada linha do arquivo associa um e-mail a uma chave. O \arq{namespaces="git"} limita a chave a assinaturas do Git.

\subsection{Passo 5 — Testar}

\begin{bashcode}
git commit --allow-empty -m "Testando a assinatura SSH"
git log --show-signature -1
\end{bashcode}
\begin{saida}
commit 86613eb51dee51c80a61f2cf8f1532df427e60c5
Good "git" signature for maria@exemplo.com with ED25519 key SHA256:Z/dW7mOFmUlo/Po1CSObYXoO3bAPY2bTohQtnJY0PGc
Author: Maria Exemplo <maria@exemplo.com>
\end{saida}

% -----------------------------------------------------------------------------
\section{Conferindo no GitHub}

Envie um commit assinado com \comando{git push} e abra a lista de commits:

\begin{itemize}
  \item Na aba \textbf{Commits} do repositório, o selo \textbf{Verified} aparece ao lado do commit. Clique nele para ver qual chave assinou.
\end{itemize}

No terminal, uma letra resume o estado de cada commit:

\begin{bashcode}
git log --format='%h %G? %s'
\end{bashcode}

\begin{center}
\renewcommand{\arraystretch}{1.3}
\begin{tabularx}{0.9\linewidth}{@{}>{\ttfamily\bfseries\color{roxoescuro}}l X@{}}
\toprule
G & assinatura válida \\
N & sem assinatura \\
B & assinatura inválida \\
E & não foi possível conferir (por exemplo, falta a chave pública ou o \arq{allowed\_signers}) \\
\bottomrule
\end{tabularx}
\end{center}

\begin{dica}[Vigilant mode no GitHub]
Em \textbf{Settings \(\rightarrow\)\ SSH and GPG keys}, a opção \textbf{Vigilant mode} (\enquote{Flag unsigned commits as unverified}) faz o GitHub marcar como \textbf{Unverified} todo commit seu que não esteja assinado, inclusive os que alguém fizer se passando por você. Ligue só depois que todos os seus computadores estiverem assinando.
\end{dica}

% -----------------------------------------------------------------------------
\section{Cuidados com a chave}

\begin{itemize}
  \item \textbf{Guarde a chave privada e a senha.} Se perder a chave, é só criar outra e cadastrar, mas perde-se a continuidade.
  \item \textbf{Chave vazou? Remova-a da conta e cadastre outra.} No GitHub, a verificação é registrada no momento do push: os commits antigos \textbf{continuam Verified} mesmo que a chave seja removida, expire ou seja revogada.
  \item \textbf{Commits feitos pela interface web} já saem assinados pela própria plataforma.
  \item No GitHub, um Pull Request integrado com \textbf{Rebase and merge} reescreve os commits, e eles perdem a sua assinatura.
  \item \textbf{Não é preciso publicar a chave em servidor de chaves} (\emph{keyserver}). Alguns tutoriais pedem isso, mas o GitHub só usa a chave cadastrada na conta.
\end{itemize}

\section{Problemas comuns}

\begin{center}
\renewcommand{\arraystretch}{1.4}
\small
\begin{tabularx}{\linewidth}{@{}>{\raggedright\arraybackslash\ttfamily\footnotesize}p{5.6cm} >{\raggedright\arraybackslash}X@{}}
\toprule
\textrm{\textbf{Mensagem ou sintoma}} & \textbf{O que fazer} \\
\midrule
error: gpg failed to sign the data & quase sempre é o \arq{GPG\_TTY} (passo 7 do GPG). Confira também se o \arq{user.signingkey} bate com \comando{gpg --list-secret-keys} \\
gpg: signing failed: Inappropriate ioctl for device & falta o \arq{export GPG\_TTY=\$(tty)} \\
No secret key / secret key not available & o ID configurado não existe neste chaveiro. No Windows, confira se o Git usa o mesmo GPG em que a chave foi criada \\
gpg.ssh.allowedSignersFile needs to be configured and exist for ssh signature verification & crie o \arq{allowed\_signers} (passo 4 do SSH). A assinatura em si está certa \\
either user.signingkey or gpg.ssh.defaultKeyCommand needs to be configured & falta o \comando{git config --global user.signingkey} \\
\textrm{Selo \textbf{Unverified} no site} & o e-mail do commit não está verificado na conta, a chave não foi cadastrada, ou a chave SSH foi cadastrada só como \emph{Authentication Key} \\
\bottomrule
\end{tabularx}
\end{center}

\begin{dica}[Fontes oficiais]
Os passos deste capítulo seguem a documentação do GitHub \cite{github_assinatura}. Os nomes dos botões podem mudar com novas versões do site.
\end{dica}
